% ================================================================
\section{Divisão e conquista}\label{sec:dc}
% ================================================================
\textbf{Paradigma:} (1) \emph{Divisão}: quebrar o problema $\Pi$ em subproblemas menores \textbf{do mesmo tipo} $\Pi_1,\dots,\Pi_k$; (2) \emph{Conquista}: resolvê-los recursivamente (casos pequenos diretamente); (3) \emph{Combinação}: juntar as soluções parciais com um procedimento ``simples''. Busca-se \textbf{bom balanceamento} (partes de tamanhos parecidos). Funciona bem quando os subproblemas são \textbf{disjuntos}; quando eles se repetem muito, a técnica certa é programação dinâmica (§\ref{sec:pd}).
Exemplos do curso: mergesort, quicksort, BFPRT, contagem de inversões, Strassen, subvetor máximo.

\begin{prova}[Q5: modelo de resposta para ``escreva o pseudocódigo por divisão e conquista'']
\begin{enumerate}
  \item \textbf{Assinatura}: \textsc{Resolve}$(V,ini,fim)$ trabalhando num intervalo (evita copiar vetores). Diga o que a função \textbf{retorna} (às vezes mais de um valor: soma \emph{e} índices; contagem \emph{e} vetor ordenado).
  \item \textbf{Caso base}: $ini=fim$ (ou $n\le$ constante): resposta direta.
  \item \textbf{Divisão}: $meio:=\floor{(ini+fim)/2}$.
  \item \textbf{Conquista}: duas chamadas recursivas, $[ini,meio]$ e $[meio+1,fim]$.
  \item \textbf{Combinação}: o que ``atravessa'' o meio (a parte criativa); geralmente $O(n)$.
  \item \textbf{Complexidade}: escreva a recorrência e resolva (árvore de recursão ou Teorema Mestre, §\ref{sec:recorrencias}): $2T(n/2)+\Theta(n)=\Theta(n\log n)$.
  \item \textbf{Correção} (uma frase): toda solução está inteira à esquerda, inteira à direita ou cruza o meio; cada caso é tratado.
\end{enumerate}
\end{prova}

% ------------------------------------------------------------------
\subsection{Contagem de inversões (aula de 11/09)}
Dada $a_1,\dots,a_n$, um par $i<j$ é uma \textbf{inversão} se $a_i>a_j$ (métrica de similaridade entre rankings: minha ordem $1,\dots,n$ vs.\ a sua). Força bruta: $\Theta(n^2)$.
\textbf{D\&C:} divide em metades $A$ e $B$; conta recursivamente as inversões dentro de cada uma; \textbf{combina} contando os pares $(a,b)$ com $a\in A$, $b\in B$, $a>b$. Se $A$ e $B$ estão \emph{ordenadas}, isso é fácil durante a intercalação: ao comparar $a_i$ e $b_j$, se $a_i\le b_j$, $a_i$ não é invertido com nenhum elemento restante de $B$; se $a_i>b_j$, então $b_j$ é invertido com \textbf{todos os elementos restantes de $A$}. Por isso a função também devolve o vetor ordenado (é um mergesort ``instrumentado'').
\begin{pseudo}{Contagem de inversões (aula + completado)}
\begin{algorithmic}
\Function{CalculaInversões}{$V,ini,fim$} \Comment{conta e deixa $V[ini..fim]$ ordenado}
  \If{$ini=fim$} \State \Return $0$ \EndIf
  \State $meio:=\floor{(ini+fim)/2}$
  \State $a:=$ \Call{CalculaInversões}{$V,ini,meio$}
  \State $b:=$ \Call{CalculaInversões}{$V,meio+1,fim$}
  \State $c:=$ \Call{MergeCalculandoInversões}{$V,ini,meio,fim$}
  \State \Return $a+b+c$
\EndFunction
\Function{MergeCalculandoInversões}{$V,ini,meio,fim$}
  \State $i:=ini$;\ $j:=meio+1$;\ $k:=1$;\ $cont:=0$
  \While{$i\le meio$ \textbf{e} $j\le fim$}
    \If{$V[i]\le V[j]$} \State $W[k]:=V[i]$;\ $i:=i+1$
    \Else
      \State $W[k]:=V[j]$;\ $j:=j+1$
      \State $cont:=cont+(meio-i+1)$ \Comment{$V[j]$ é menor que todos os restantes da esquerda}
    \EndIf
    \State $k:=k+1$
  \EndWhile
  \State copie os restantes de $V[i..meio]$ e de $V[j..fim]$ para $W$;\ copie $W$ para $V[ini..fim]$
  \State \Return $cont$
\EndFunction
\end{algorithmic}
\end{pseudo}
$T(n)=2T(n/2)+\Theta(n)\Rightarrow\Theta(n\log n)$. \emph{Cuidado:} usar ``$\le$'' (iguais não são inversão) e somar $meio-i+1$ (quantos restam em $A$), não 1. Ex.\ conferido: $1,5,4,8,10,2,6,9,12,11,3,7$ tem \textbf{22} inversões.

% ------------------------------------------------------------------
\subsection{Subvetor contíguo de soma máxima (Teste 2, Exercício 1)}
Dado $A[1..n]$ (inteiros, podem ser negativos), achar $i\le j$ maximizando $A[i]+\dots+A[j]$ (subvetor \textbf{não vazio}).
\textbf{Ideia:} o melhor subvetor de $A[ini..fim]$ está (i) inteiro na metade esquerda, (ii) inteiro na direita, ou (iii) \textbf{cruza o meio}, i.e., é $A[i..meio]\cup A[meio+1..j]$. O caso (iii) se resolve em tempo linear: melhor sufixo da esquerda terminando em $meio$ + melhor prefixo da direita começando em $meio+1$.
\begin{pseudo}{Subvetor de soma máxima por D\&C (retorna soma, início, fim)}
\begin{algorithmic}
\Function{MaxSub}{$A,ini,fim$}
  \If{$ini=fim$} \State \Return $(A[ini],\,ini,\,ini)$ \Comment{não vazio: um elemento}
  \EndIf
  \State $meio:=\floor{(ini+fim)/2}$
  \State $(s_E,i_E,j_E):=$ \Call{MaxSub}{$A,ini,meio$}
  \State $(s_D,i_D,j_D):=$ \Call{MaxSub}{$A,meio+1,fim$}
  \State $(s_C,i_C,j_C):=$ \Call{MaxCruzado}{$A,ini,meio,fim$}
  \State \Return a tripla de maior soma entre $(s_E,i_E,j_E)$, $(s_D,i_D,j_D)$, $(s_C,i_C,j_C)$
\EndFunction
\Function{MaxCruzado}{$A,ini,meio,fim$}
  \State $soma:=0$;\ $melhorE:=-\infty$
  \For{$i:=meio$ \textbf{decrescendo até} $ini$}
    \State $soma:=soma+A[i]$
    \If{$soma>melhorE$} \State $melhorE:=soma$;\ $iMax:=i$ \EndIf
  \EndFor
  \State $soma:=0$;\ $melhorD:=-\infty$
  \For{$j:=meio+1$ \textbf{até} $fim$}
    \State $soma:=soma+A[j]$
    \If{$soma>melhorD$} \State $melhorD:=soma$;\ $jMax:=j$ \EndIf
  \EndFor
  \State \Return $(melhorE+melhorD,\ iMax,\ jMax)$
\EndFunction
\end{algorithmic}
\end{pseudo}
\begin{itemize}
  \item \textbf{Complexidade:} $T(n)=2T(n/2)+\Theta(n)\Rightarrow\Theta(n\log n)$ (Cormen §4.1).
  \item \textbf{Todos negativos:} funciona, pois o caso base devolve um elemento e o cruzado usa pelo menos $A[meio]$ e $A[meio+1]$; o máximo das três opções é o maior elemento. Ex.\ 2 do teste: $[-8,-3,-6,-2,-5,-4]\to -2$.
  \item Ex.\ 1 do teste (conferido): $[-2,1,-3,4,-1,2,1,-5,4]\to$ soma $6$, subvetor $[4,-1,2,1]$. \emph{O enunciado indexa a partir de 0} (posições 3 a 6); com índices a partir de 1, posições 4 a 7.
  \item (Fora de D\&C: o algoritmo de Kadane resolve em $\Theta(n)$ mantendo o melhor sufixo terminando em cada posição. Também existe D\&C linear devolvendo (soma total, melhor prefixo, melhor sufixo, melhor) de cada metade: $T(n)=2T(n/2)+O(1)=\Theta(n)$.)
\end{itemize}

% ------------------------------------------------------------------
\subsection{Multiplicação de matrizes: Strassen (Teste 2, Exercício 2)}
Produto por blocos ingênuo: 8 multiplicações de $n/2\times n/2$, $T(n)=8T(n/2)+\Theta(n^2)=\Theta(n^3)$. Strassen usa \textbf{7}: $T(n)=7T(n/2)+\Theta(n^2)=\Theta(n^{\log_27})\approx\Theta(n^{2{,}807})$ (Teorema Mestre caso 1, pois $n^2=O(n^{\log_27-\varepsilon})$).
\begin{pseudo}{Strassen$(A,B,n)$ -- $n$ potência de 2; Soma/Subtrai são dadas}
\begin{algorithmic}
\Function{Strassen}{$A,B,n$}
  \If{$n=1$} \State \Return $[\,A[1,1]\cdot B[1,1]\,]$ \EndIf
  \State particione $A=\begin{bmatrix}A_{11}&A_{12}\\A_{21}&A_{22}\end{bmatrix}$, $B=\begin{bmatrix}B_{11}&B_{12}\\B_{21}&B_{22}\end{bmatrix}$ em blocos $\frac n2\times\frac n2$
  \State $M_1:=$ \Call{Strassen}{$A_{11}+A_{22},\ B_{11}+B_{22},\ n/2$}
  \State $M_2:=$ \Call{Strassen}{$A_{21}+A_{22},\ B_{11},\ n/2$}
  \State $M_3:=$ \Call{Strassen}{$A_{11},\ B_{12}-B_{22},\ n/2$}
  \State $M_4:=$ \Call{Strassen}{$A_{22},\ B_{21}-B_{11},\ n/2$}
  \State $M_5:=$ \Call{Strassen}{$A_{11}+A_{12},\ B_{22},\ n/2$}
  \State $M_6:=$ \Call{Strassen}{$A_{21}-A_{11},\ B_{11}+B_{12},\ n/2$}
  \State $M_7:=$ \Call{Strassen}{$A_{12}-A_{22},\ B_{21}+B_{22},\ n/2$}
  \State $C_{11}:=M_1+M_4-M_5+M_7$;\quad $C_{12}:=M_3+M_5$
  \State $C_{21}:=M_2+M_4$;\quad $C_{22}:=M_1-M_2+M_3+M_6$
  \State \Return $\begin{bmatrix}C_{11}&C_{12}\\C_{21}&C_{22}\end{bmatrix}$
\EndFunction
\end{algorithmic}
\end{pseudo}
São 18 somas/subtrações de matrizes $\frac n2\times\frac n2$ ($\Theta(n^2)$) por chamada. Na prática usa-se o método clássico abaixo de um certo $n$ (constantes maiores, menos estabilidade numérica).

% ------------------------------------------------------------------
\subsection{Outros exemplos clássicos de D\&C (treino para a Q5)}
\begin{pseudo}{Busca binária, potência rápida e máximo/mínimo}
\begin{algorithmic}
\Function{BuscaBin}{$V,ini,fim,x$} \Comment{$V$ ordenado; $T(n)=T(n/2)+O(1)=O(\log n)$}
  \If{$ini>fim$} \State \Return $-1$ \EndIf
  \State $meio:=\floor{(ini+fim)/2}$
  \If{$V[meio]=x$} \State \Return $meio$
  \ElsIf{$x<V[meio]$} \State \Return \Call{BuscaBin}{$V,ini,meio-1,x$}
  \Else \State \Return \Call{BuscaBin}{$V,meio+1,fim,x$}
  \EndIf
\EndFunction
\Function{Pot}{$a,n$} \Comment{$a^n$ com $O(\log n)$ multiplicações}
  \If{$n=0$} \State \Return $1$ \EndIf
  \State $t:=$ \Call{Pot}{$a,\floor{n/2}$}
  \If{$n$ par} \State \Return $t\cdot t$ \Else \State \Return $t\cdot t\cdot a$ \EndIf
\EndFunction
\Function{MaxMin}{$V,ini,fim$} \Comment{$\approx 3n/2$ comparações (contra $2n-2$ do ingênuo)}
  \If{$ini=fim$} \State \Return $(V[ini],V[ini])$ \EndIf
  \If{$fim=ini+1$} \State \Return $(\max,\min)$ de $V[ini],V[fim]$ com 1 comparação \EndIf
  \State $meio:=\floor{(ini+fim)/2}$
  \State $(M_1,m_1):=$ \Call{MaxMin}{$V,ini,meio$};\ $(M_2,m_2):=$ \Call{MaxMin}{$V,meio+1,fim$}
  \State \Return $(\max\{M_1,M_2\},\ \min\{m_1,m_2\})$
\EndFunction
\end{algorithmic}
\end{pseudo}

\begin{pseudo}{Elemento majoritário (aparece $>n/2$ vezes), $T(n)=2T(n/2)+\Theta(n)=\Theta(n\log n)$}
\begin{algorithmic}
\Function{Maj}{$V,ini,fim$} \Comment{retorna o majoritário de $V[ini..fim]$ ou \textsc{nenhum}}
  \If{$ini=fim$} \State \Return $V[ini]$ \EndIf
  \State $meio:=\floor{(ini+fim)/2}$
  \State $a:=$ \Call{Maj}{$V,ini,meio$};\ $b:=$ \Call{Maj}{$V,meio+1,fim$}
  \For{$c\in\{a,b\}$, $c\ne$ \textsc{nenhum}}
    \If{$c$ ocorre mais de $(fim-ini+1)/2$ vezes em $V[ini..fim]$} \State \Return $c$ \EndIf \Comment{contagem $O(n)$}
  \EndFor
  \State \Return \textsc{nenhum}
\EndFunction
\end{algorithmic}
\end{pseudo}
\emph{Correção:} se $x$ é majoritário no todo, é majoritário em pelo menos uma das metades (senão ocorreria $\le n/2$ vezes).

\paragraph{Mais ideias de D\&C.} \emph{Karatsuba} (inteiros de $n$ dígitos com 3 produtos de $n/2$: $\Theta(n^{\log_23})\approx n^{1{,}585}$); \emph{par de pontos mais próximos} no plano ($\Theta(n\log n)$: divide pela mediana de $x$, combina olhando a faixa central com no máximo 7 vizinhos por ponto); \emph{pico em vetor unimodal} e \emph{raiz de função monótona} (busca binária, $O(\log n)$); \emph{contar ocorrências de $x$ em vetor ordenado} (duas buscas binárias).

\begin{cuidado}[Erros frequentes em pseudocódigo de D\&C]
\begin{itemize}
  \item Esquecer o caso base ou fazê-lo errado (ex.: subvetor \emph{vazio} permitido quando o enunciado pede não vazio).
  \item Dividir em $[ini,meio-1]$ e $[meio,fim]$ com $meio=\floor{(ini+fim)/2}$: com 2 elementos gera recursão infinita. Use $[ini,meio]$ e $[meio+1,fim]$.
  \item Combinação que olha só um elemento de cada lado (o subvetor cruzado precisa do \emph{melhor sufixo} e do \emph{melhor prefixo}).
  \item Não devolver a informação extra necessária (índices; vetor ordenado nas inversões).
  \item Não dizer a complexidade (o enunciado sempre pede!).
\end{itemize}
\end{cuidado}

\begin{chave}
D\&C = dividir (meio), conquistar (2 chamadas), combinar (o que cruza o meio). $2T(n/2)+\Theta(n)=\Theta(n\log n)$; $T(n/2)+O(1)=\Theta(\log n)$; $7T(n/2)+\Theta(n^2)=\Theta(n^{2{,}807})$; $T(n/5)+T(3n/4)+\Theta(n)=\Theta(n)$.
\end{chave}
